\section{The Concrete Syntax and the Constraints}
\label{sec:constraints}

The metamodels of Section~\ref{sec:sourcemetamodel} and
Section~\ref{sec:targetmetamodel} define what can be said. Two things are needed
before a model can be written down and trusted: a concrete syntax an author
types, and the constraints that refuse a model the abstract syntax alone would
admit.

\subsection{The Concrete Syntax}

The source language has a textual concrete syntax, defined by an Xtext grammar
over the metamodel rather than the other way round
(Section~\ref{sec:alternatives}). There are two grammars, one per authored
document, and both import the registered metamodel package and return its
classes: a rule never declares a type of its own, so no concept can enter the
language through the syntax without existing in the metamodel first.

The syntax is the YAML subset the estate's files already use, which keeps the
authored documents readable by the tools that read them today. Three properties
of the grammar are worth reporting, because each is a consequence of that
choice.

\textbf{Indentation is tokenised, not parsed.} Block structure reaches the
parser as synthetic \texttt{BEGIN} and \texttt{END} tokens produced ahead of it,
so an indented block, a flow mapping and a flow sequence all parse through one
rule rather than three.

\textbf{Keys are unordered.} Every rule is an unordered group, because the order
of keys in a mapping carries no meaning. Writing the rules as sequences would
make the grammar refuse documents the model accepts.

\textbf{Names are references.} Where a route names the process and the surface
it serves, the grammar declares a cross-reference rather than a string, so the
name is resolved by the linker against the declared objects. A name with no
target is a link that does not resolve, which is why a scrape naming an
undeclared process (\texttt{E\_UNKNOWN\_PROCESS}) and a tier naming an
undeclared proxy (\texttt{E\_UNKNOWN\_TIER\_PROXY}) are not invariants in
Table~\ref{tab:constraints}: the linker has already found them.

The editor is the generated Xtext one, with the constraints below loaded into
it: the Complete OCL document binds to the same registered package the parser
instantiates, so an invariant is evaluated against the objects the editor has
just built, and a violation is reported in the editor rather than at render
time.

\subsection{The Constraints}

Fifteen invariants are written in Complete OCL against the source metamodel,
carrying thirteen distinct diagnostic codes: two codes each sit on two contexts,
because the same defect is reachable from two classes.

The document follows one naming rule, which is what makes a diagnostic
self-locating: \textbf{an invariant is named by the code it emits, and its
context is the object the diagnostic points at.} A violation is therefore a
\texttt{(code, path)} pair with no lookup table in between, where the path is
the RFC 6901 JSON Pointer of the refused object.

\begin{table}[!htbp]
\centering
\footnotesize
\begin{tabular}{>{\raggedright\arraybackslash}p{0.35\linewidth}>{\raggedright\arraybackslash}p{0.14\linewidth}>{\raggedright\arraybackslash}p{0.43\linewidth}}
\hline
\textbf{Code} & \textbf{Context} & \textbf{What it refuses} \\
\hline
E\_ALERT\_CLASS\_WITHOUT\_SIGNAL & Observability & A class stating how loudly to wake someone, with no signal to wake them about. \\
E\_CUTOVER\_UNHONOURABLE & Process & A rolling cutover over storage that cannot attach twice, which no rollout can honour. \\
E\_ENGINE\_WITHOUT\_DURABILITY & Process & A backup method named where no volume's durability class asks for one. \\
E\_DURABILITY\_WITHOUT\_ENGINE & Volume & A volume asking for a backup whose method would have to be guessed. \\
E\_ENV\_CANNOT\_RELOAD & Rotation & A secret that tolerates a reload, delivered as an environment variable, which is read once at start. \\
E\_ILLEGAL\_DELIVERY\_FOR\_ACCESS & KvGrant & A delivery the access tier has nothing to project at render time. \\
E\_NON\_KV\_DELIVERY & TransitGrant, DatabaseGrant & A grant delivered as a value, where the credential is used or minted rather than read. \\
E\_DUPLICATE\_ROUTE\_MATCH & Route & Two routes of one exposure sharing a path and a match, which derived precedence cannot order. \\
E\_NO\_FORWARD\_AUTH\_ENDPOINT & Tier & A tier serving authenticated visitors with nothing named to decide whether one is allowed in. \\
E\_NO\_TIER\_FOR\_AUDIENCE & Exposure, Route & An audience no declared edge tier serves, so the route reaches no edge. \\
E\_NO\_ENGINE\_POLICY & Process & An engine the platform declares no backup method for. \\
E\_NO\_DURABILITY\_POLICY & Volume & A durability class the platform declares no policy for. \\
E\_SECRETS\_AT\_REST\_REQUIRED & Grant & A secret written into a cluster that does not encrypt secrets at rest. \\
\hline
\end{tabular}
\caption{The Constraints, Their Context and What Each Refuses}
\label{tab:constraints}
\end{table}

The last four read facts from the Platform document while sitting on a class of
the project document. That placement is deliberate: an invariant sits on the
object whose author must change something, and it is a project file, not the
platform, that has to be edited when a volume asks for a guarantee the estate
does not keep. They reach the other document through \texttt{allInstances()},
whose extent is the resource set the document was read into, and each is guarded
so that it holds trivially when a project file is read on its own:

\begin{lstlisting}[basicstyle=\footnotesize\ttfamily, frame=single, captionpos=b,
                   caption={E\_NO\_TIER\_FOR\_AUDIENCE, guarded so a project file read alone still validates},
                   label={lst:notier}]
context Exposure
    inv E_NO_TIER_FOR_AUDIENCE:
        Tier.allInstances()->notEmpty() implies
            Tier.allInstances()->exists(audiences->includes(self.audience))
\end{lstlisting}

One modelling detail shows up in several invariants and is worth naming, because
it is a property of EMF rather than of the domain. EMF reads an unset
enumeration attribute as its first literal, so an optional enumeration cannot be
distinguished from one written with its first value. Each such vocabulary
therefore carries a literal of its own for absence, \texttt{absent}, which a
document cannot write, and the invariants test against it rather than against
\texttt{null}.

Constraints that need every project document at once, such as two applications
claiming one identifier, are not in this document. They belong to the
composition step, where the union of the authored documents exists, and they are
checked there against the fixtures in the negative set.

\subsection{Every Constraint Has a Fixture}

A constraint that has only ever been satisfied is not evidence that it fires.
Each code in Table~\ref{tab:constraints} is therefore paired with a committed
model that violates it, and with a committed oracle file naming the set of
\texttt{(code, path)} pairs that model must produce. The four cross-document
codes are fixtures of two files, a project document and the platform document it
contradicts, because one file alone cannot express the defect.

The pairing is checked rather than asserted: each refused model is run through
the parser and validated, and the diagnostics it produces are compared with the
oracle. The accepted examples of Section~\ref{sec:examplemodels} prove the
language can express the estate; these prove it refuses what it says it refuses.

\subsection{The Second Implementation}

The same source language is declared a second time, by hand, in the TypeScript
production implementation this repository exists to build. Neither declaration
is generated from the other, and neither is the oracle for the other: each is
tested on its own against the committed oracle files, and agreement between the
two follows from both agreeing with the oracle.

They are compared on the four things the project proposal names: validation, so
the parsed model of each accepted example matches one committed canonical JSON
document; dependency resolution, so the resolved edges match a second; errors,
so each refused model produces the same set of \texttt{(code, document, path)}
triples; and generated resources, so the generated tree matches the committed
one byte for byte. They are required to agree on nothing else, and in one
respect they deliberately differ: the production implementation keeps the
resolved model and the typed target objects apart, while the model-driven
implementation holds both in the one target metamodel of
Section~\ref{sec:targetmetamodel}, because a model-to-text template reads one
model.

Two independent declarations of one language, held to shared oracles, is also
what makes the metamodel falsifiable in a way a single implementation cannot be:
a concept that only one of them can express is a concept the model has not
actually pinned down.
